% --- Содержимое файла materials/4_1.tex ---

\label{seventh-bilet}
\subsection{Теорема Рисса}
\begin{definition}[Линейное нормированное пространство]
	Линейным нормированным пространством над полем $\mathbb{K}$ ($\mathbb{R}$ или $\mathbb{C}$) называется линейное пространство $E$ с функцией $\|\cdot\| : E \to \mathbb{R}$, обладающей свойствами:
	\begin{enumerate}
		\item $\forall x \in E: \|x\| \geqslant 0$, причем $\|x\| = 0 \Leftrightarrow x = 0$
		\item $\forall x \in E: \forall \alpha \in \mathbb{K}: \|\alpha x\| = |\alpha|\|x\|$
		\item $\forall x, y \in E: \|x + y\| \leqslant \|x\| + \|y\|$ (неравенство треугольника)
	\end{enumerate}
	Функция $\|\cdot\|$ называется \textbf{нормой} на пространстве $E$.
\end{definition}

\begin{example}
	Любое линейное нормированное пространство $E$ является метрическим с метрикой $\rho(x, y) := \|x - y\|$. Приведенные ранее метрические пространства являются линейными нормированными, и метрика в них имеет вид нормы разности элементов.
\end{example}

\begin{remark}
	Норма непрерывна как функция $E \to \mathbb{R}$. Если $x_n \xrightarrow[E]{} x$, то $\rho(x_n, x) = \|x_n - x\| \to 0$. По неравенству треугольника:
	\begin{itemize}
		\item $\|x_n\| \leqslant \|x_n - x\| + \|x\|$
		\item $\|x\| \leqslant \|x_n - x\| + \|x_n\|$
	\end{itemize}
	Таким образом, $\left|\|x_n\| - \|x\|\right| \to 0$, то есть $\|x_n\| \to \|x\|$.
\end{remark}

\begin{definition}
	Пусть $E$ --- линейное нормированное пространство. Множество $L \subset E$ называется:
	\begin{itemize}
		\item \textbf{Линейным многообразием}, если $L$ замкнуто относительно сложения и умножения на скаляры.
		\item \textbf{Подпространством}, если $L$ является линейным многообразием и при этом замкнуто как множество.
	\end{itemize}
\end{definition}

\begin{definition}[Линейная оболочка]
	Линейной оболочкой множества $M \subset E$ называется множество вида:
	\[[M] := \left\{\sum_{k = 1}^n \alpha_i m_i : \alpha_i \in \mathbb{K}, m_i \in M\right\}\]
\end{definition}

\begin{remark}
	Линейная оболочка всегда является линейным многообразием, но не всегда является подпространством.
\end{remark}

\begin{example}
	Пусть $E = C[0,1]$ и $M$ — многочлены. Тогда $M = [M]$ — многообразие, но не подпространство. По теореме Вейерштрасса, предел последовательности многочленов может не быть многочленом, значит $[M]$ не содержит все свои точки прикосновения (не замкнуто).
\end{example}

\begin{definition}[Эквивалентность норм]
	Нормы $\|\cdot\|$ и $\|\cdot\|'$ на $E$ называются \textbf{эквивалентными}, если $\exists C_1, C_2 > 0$ такие, что для любого $x \in E$:
	\[C_1\|x\|' \leqslant \|x\| \leqslant C_2\|x\|'\]
\end{definition}

\begin{definition}
	\textbf{Размерностью} пространства $E$ ($\dim E$) называется наибольший размер линейно независимой системы в $E$.
\end{definition}

\begin{proposition}[Лемма о «почти перпендикуляре»]\label{almostperp}
	Пусть $E$ --- ЛНП, $L \subsetneq E$ --- подпространство. Тогда для любого $\varepsilon > 0$ существует $y \in E$ такой, что $\|y\| = 1$ и $\rho(y, L) = \inf_{z \in L}\|y - z\| \geqslant 1 - \varepsilon$.
\end{proposition}

\begin{proof}
	\begin{enumerate}
		\item Зафиксируем $y_0 \in E \setminus L$ и положим $d := \rho(y_0, L) > 0$.
		\item Выберем $z_0 \in L$ такой, что $d \leqslant \|y_0 - z_0\| \leqslant \frac{d}{1-\varepsilon}$. 
		\item Покажем, что подходит вектор $y := \frac{y_0 - z_0}{\|y_0 - z_0\|}$. Очевидно, $\|y\| = 1$.
		\item Положим $\alpha = \frac{1}{\|y_0 - z_0\|}$. Для любого $z \in L$:
		\[ \|y - z\| = \|\alpha(y_0 - z_0) - z\| = \alpha \| y_0 - (z_0 + \frac{1}{\alpha} z) \| \]
		\item Так как $(z_0 + \frac{1}{\alpha} z) \in L$, то $\|y_0 - (z_0 + \frac{1}{\alpha} z)\| \geqslant d$. Тогда:
		\[ \|y - z\| \geqslant \alpha d = \frac{d}{\|y_0 - z_0\|} \geqslant \frac{d}{d/(1-\varepsilon)} = 1 - \varepsilon \]
		Значит, $\inf_{z \in L}\|y - z\| \geqslant 1 - \varepsilon$.
	\end{enumerate}
\end{proof}

\begin{idea}[Лемма о почти перпендикуляре]
	\textbf{Геометрический смысл:} В конечномерном случае мы всегда можем найти вектор, строго перпендикулярный плоскости. В произвольном ЛНП «идеального» перпендикуляра может не существовать, но лемма утверждает, что мы всегда можем найти вектор $y$, который «почти» перпендикулярен подпространству $L$.
	
	\textbf{Важное уточнение по расстоянию:}
	\begin{itemize}
		\item Мы измеряем расстояние от \textbf{конца (носка)} вектора $y$ до подпространства $L$. 
		\item Хотя «начало» вектора (пятка) находится в нуле и принадлежит $L$, сам вектор «торчит» в сторону сферы. 
		\item Так как длина вектора равна $1$, кратчайший путь от его конца до $L$ \textbf{не может быть больше 1} (путь в ноль всегда доступен).
		\item Лемма гарантирует, что мы можем направить вектор так, чтобы его конец был удален от $L$ \textbf{почти на всю свою длину} ($\geqslant 1 - \varepsilon$).
	\end{itemize}
	
	\textbf{Суть доказательства:} Мы берем вектор $y_0$ вне подпространства и «нормируем» его отклонение от почти ближайшей точки $z_0 \in L$. Чем ближе $z_0$ к истинному инфимуму, тем «вертикальнее» стоит итоговый вектор $y$ относительно $L$.
\end{idea}

\begin{theorem}[Теорема Рисса]\label{thm4.1}
	Пусть $E$ --- ЛНП. Тогда единичная сфера $S(0, 1)$ компактна в $E \Leftrightarrow \dim E < +\infty$.
\end{theorem}

\begin{proof}
	\begin{itemize}
		\item[$\Leftarrow$] В конечномерном пространстве все нормы эквивалентны. Относительно евклидовой нормы сфера $S(0, 1)$ компактна по теореме Гейне-Бореля.
		
		\item[$\Rightarrow$] \textbf{От противного.} Пусть $\dim E = \infty$. Построим последовательность $\{x_n\} \subset S(0, 1)$ с расстояниями $\rho(x_i, x_j) \geqslant 1 - \varepsilon$:
		\begin{enumerate}
			\item Выберем $x_1 \in S(0, 1)$ произвольно.
			\item Пусть $L_1 = [x_1]$. Т.к. $\dim L_1 < \infty$, оно замкнуто и является подпространством. По лемме~\ref{almostperp} выберем $x_2 \in S(0, 1)$ так, что $\rho(x_2, L_1) \geqslant 1 - \varepsilon$.
			\item Пусть $L_2 = [x_1, x_2]$. Оно также замкнуто. Выберем $x_3 \in S(0, 1)$ так, что $\rho(x_3, L_2) \geqslant 1 - \varepsilon$.
			\item Процесс не закончится, так как $\dim E = \infty$. Получим последовательность, из которой нельзя выделить сходящуюся (т.к. $\rho(x_n, x_m) \geqslant 1-\varepsilon$). Значит, сфера не вполне ограничена и не компактна. Противоречие.
		\end{enumerate}
	\end{itemize}
\end{proof}	

\begin{idea}[Смысл теоремы Рисса]
	\textbf{Суть:} Это критерий того, насколько «просторно» в пространстве. 
	
	\begin{itemize}
		\item \textbf{В конечномерном мире} (как $\mathbb{R}^3$) сфера «тесная». Если наставить на ней бесконечно много точек, им придется прижиматься друг к другу (сгущаться). Это обеспечивает компактность.
		\item \textbf{В бесконечномерном мире} всегда есть «куда расти». Там можно расставить бесконечный «ёж» из векторов так, что все они будут на расстоянии 1 друг от друга и никогда не сблизятся. 
	\end{itemize}
	\textbf{Вывод:} Сфера компактна только тогда, когда пространство конечно. Бесконечномерность «убивает» компактность из-за бесконечного избытка места.
\end{idea}

\begin{idea}[Теорема Рисса]
	\textbf{Почему это важно:} Эта теорема проводит жесткую границу между привычным нам миром (конечномерным) и миром бесконечных функций. 
	
	\textbf{Геометрия бесконечномерности:}
	\begin{itemize}
		\item В конечномерном пространстве «место» ограничено. Если мы пытаемся расставить бесконечно много точек на сфере, им приходится сгущаться, образуя предельную точку.
		\item В бесконечномерном пространстве появляется «бесконечно много новых направлений». Благодаря лемме о почти перпендикуляре, мы можем каждый раз выбирать вектор, который \textbf{почти перпендикулярен всем предыдущим}.
		\item В итоге мы строим последовательность $\{x_n\}$, где расстояние между любыми двумя точками $\rho(x_n, x_m) \approx 1$. 
	\end{itemize}
	
	\textbf{Суть противоречия:} 
	Такая последовательность «разлетается» во все стороны. Она никогда не сможет сблизиться (стать фундаментальной), а значит, из нее нельзя выделить сходящуюся подпоследовательность. Сфера оказывается «слишком просторной», чтобы быть компактной.
\end{idea}